% Options for packages loaded elsewhere
\PassOptionsToPackage{unicode}{hyperref}
\PassOptionsToPackage{hyphens}{url}
\PassOptionsToPackage{dvipsnames,svgnames,x11names}{xcolor}
\documentclass[
  12pt,
]{article}
\usepackage{xcolor}
\usepackage[margin=1in]{geometry}
\usepackage{amsmath,amssymb}
\setcounter{secnumdepth}{5}
\usepackage{iftex}
\ifPDFTeX
  \usepackage[T1]{fontenc}
  \usepackage[utf8]{inputenc}
  \usepackage{textcomp} % provide euro and other symbols
\else % if luatex or xetex
  \usepackage{unicode-math} % this also loads fontspec
  \defaultfontfeatures{Scale=MatchLowercase}
  \defaultfontfeatures[\rmfamily]{Ligatures=TeX,Scale=1}
\fi
\usepackage{lmodern}
\ifPDFTeX\else
  % xetex/luatex font selection
  \setmainfont[]{Calibri}
\fi
% Use upquote if available, for straight quotes in verbatim environments
\IfFileExists{upquote.sty}{\usepackage{upquote}}{}
\IfFileExists{microtype.sty}{% use microtype if available
  \usepackage[]{microtype}
  \UseMicrotypeSet[protrusion]{basicmath} % disable protrusion for tt fonts
}{}
\makeatletter
\@ifundefined{KOMAClassName}{% if non-KOMA class
  \IfFileExists{parskip.sty}{%
    \usepackage{parskip}
  }{% else
    \setlength{\parindent}{0pt}
    \setlength{\parskip}{6pt plus 2pt minus 1pt}}
}{% if KOMA class
  \KOMAoptions{parskip=half}}
\makeatother
\usepackage{color}
\usepackage{fancyvrb}
\newcommand{\VerbBar}{|}
\newcommand{\VERB}{\Verb[commandchars=\\\{\}]}
\DefineVerbatimEnvironment{Highlighting}{Verbatim}{commandchars=\\\{\}}
% Add ',fontsize=\small' for more characters per line
\usepackage{framed}
\definecolor{shadecolor}{RGB}{248,248,248}
\newenvironment{Shaded}{\begin{snugshade}}{\end{snugshade}}
\newcommand{\AlertTok}[1]{\textcolor[rgb]{0.94,0.16,0.16}{#1}}
\newcommand{\AnnotationTok}[1]{\textcolor[rgb]{0.56,0.35,0.01}{\textbf{\textit{#1}}}}
\newcommand{\AttributeTok}[1]{\textcolor[rgb]{0.13,0.29,0.53}{#1}}
\newcommand{\BaseNTok}[1]{\textcolor[rgb]{0.00,0.00,0.81}{#1}}
\newcommand{\BuiltInTok}[1]{#1}
\newcommand{\CharTok}[1]{\textcolor[rgb]{0.31,0.60,0.02}{#1}}
\newcommand{\CommentTok}[1]{\textcolor[rgb]{0.56,0.35,0.01}{\textit{#1}}}
\newcommand{\CommentVarTok}[1]{\textcolor[rgb]{0.56,0.35,0.01}{\textbf{\textit{#1}}}}
\newcommand{\ConstantTok}[1]{\textcolor[rgb]{0.56,0.35,0.01}{#1}}
\newcommand{\ControlFlowTok}[1]{\textcolor[rgb]{0.13,0.29,0.53}{\textbf{#1}}}
\newcommand{\DataTypeTok}[1]{\textcolor[rgb]{0.13,0.29,0.53}{#1}}
\newcommand{\DecValTok}[1]{\textcolor[rgb]{0.00,0.00,0.81}{#1}}
\newcommand{\DocumentationTok}[1]{\textcolor[rgb]{0.56,0.35,0.01}{\textbf{\textit{#1}}}}
\newcommand{\ErrorTok}[1]{\textcolor[rgb]{0.64,0.00,0.00}{\textbf{#1}}}
\newcommand{\ExtensionTok}[1]{#1}
\newcommand{\FloatTok}[1]{\textcolor[rgb]{0.00,0.00,0.81}{#1}}
\newcommand{\FunctionTok}[1]{\textcolor[rgb]{0.13,0.29,0.53}{\textbf{#1}}}
\newcommand{\ImportTok}[1]{#1}
\newcommand{\InformationTok}[1]{\textcolor[rgb]{0.56,0.35,0.01}{\textbf{\textit{#1}}}}
\newcommand{\KeywordTok}[1]{\textcolor[rgb]{0.13,0.29,0.53}{\textbf{#1}}}
\newcommand{\NormalTok}[1]{#1}
\newcommand{\OperatorTok}[1]{\textcolor[rgb]{0.81,0.36,0.00}{\textbf{#1}}}
\newcommand{\OtherTok}[1]{\textcolor[rgb]{0.56,0.35,0.01}{#1}}
\newcommand{\PreprocessorTok}[1]{\textcolor[rgb]{0.56,0.35,0.01}{\textit{#1}}}
\newcommand{\RegionMarkerTok}[1]{#1}
\newcommand{\SpecialCharTok}[1]{\textcolor[rgb]{0.81,0.36,0.00}{\textbf{#1}}}
\newcommand{\SpecialStringTok}[1]{\textcolor[rgb]{0.31,0.60,0.02}{#1}}
\newcommand{\StringTok}[1]{\textcolor[rgb]{0.31,0.60,0.02}{#1}}
\newcommand{\VariableTok}[1]{\textcolor[rgb]{0.00,0.00,0.00}{#1}}
\newcommand{\VerbatimStringTok}[1]{\textcolor[rgb]{0.31,0.60,0.02}{#1}}
\newcommand{\WarningTok}[1]{\textcolor[rgb]{0.56,0.35,0.01}{\textbf{\textit{#1}}}}
\usepackage{longtable,booktabs,array}
\newcounter{none} % for unnumbered tables
\usepackage{calc} % for calculating minipage widths
% Correct order of tables after \paragraph or \subparagraph
\usepackage{etoolbox}
\makeatletter
\patchcmd\longtable{\par}{\if@noskipsec\mbox{}\fi\par}{}{}
\makeatother
% Allow footnotes in longtable head/foot
\IfFileExists{footnotehyper.sty}{\usepackage{footnotehyper}}{\usepackage{footnote}}
\makesavenoteenv{longtable}
\usepackage{graphicx}
\makeatletter
\newsavebox\pandoc@box
\newcommand*\pandocbounded[1]{% scales image to fit in text height/width
  \sbox\pandoc@box{#1}%
  \Gscale@div\@tempa{\textheight}{\dimexpr\ht\pandoc@box+\dp\pandoc@box\relax}%
  \Gscale@div\@tempb{\linewidth}{\wd\pandoc@box}%
  \ifdim\@tempb\p@<\@tempa\p@\let\@tempa\@tempb\fi% select the smaller of both
  \ifdim\@tempa\p@<\p@\scalebox{\@tempa}{\usebox\pandoc@box}%
  \else\usebox{\pandoc@box}%
  \fi%
}
% Set default figure placement to htbp
\def\fps@figure{htbp}
\makeatother
\setlength{\emergencystretch}{3em} % prevent overfull lines
\providecommand{\tightlist}{%
  \setlength{\itemsep}{0pt}\setlength{\parskip}{0pt}}
\usepackage{float}
\usepackage{booktabs}
\usepackage{caption}
\captionsetup{font=it, justification=centering}
\floatplacement{figure}{H}
\usepackage{bookmark}
\IfFileExists{xurl.sty}{\usepackage{xurl}}{} % add URL line breaks if available
\urlstyle{same}
% fallback for those not using the hyperref driver hyperxmp:
\makeatletter
\@ifundefined{xmpquote}{\newcommand{\xmpquote}[1]{#1}}{}
\makeatother
\hypersetup{
  pdftitle={Using Spatial Lag{,} Spatial Error{,} and Geographically Weighted Regression to Predict Median House Values in Philadelphia Block Groups},
  pdfauthor={Sara Cerami, Gautami Kasat, Tegan Lochner},
  colorlinks=true,
  linkcolor={blue},
  filecolor={Maroon},
  citecolor={Blue},
  urlcolor={Blue},
  pdfcreator={LaTeX via pandoc}}

\title{Using Spatial Lag, Spatial Error, and Geographically Weighted
Regression to Predict Median House Values in Philadelphia Block Groups}
\usepackage{etoolbox}
\makeatletter
\providecommand{\subtitle}[1]{% add subtitle to \maketitle
  \apptocmd{\@title}{\par {\large #1 \par}}{}{}
}
\makeatother
\subtitle{CPLN 671 / MUSA 5000: Statistics and Data Mining}
\author{Sara Cerami, Gautami Kasat, Tegan Lochner}
\date{November 6, 2026}

\begin{document}
\maketitle

{
\hypersetup{linkcolor=blue}
\setcounter{tocdepth}{2}
\tableofcontents
}
\newpage

\section{Introduction}\label{introduction}

In Philadelphia, housing values vary substantially. This report examines
median house values across Philadelphia using spatial statistical
methods to investigate both the relationships between housing values and
neighborhood characteristics and the extent to which these relationships
exhibit spatial dependence.

In the previous report, ordinary least squares (OLS) regression was used
to examine the relationship between median house values and five
explanatory variables (predictors): the percentage of residents with a
bachelor's degree or higher (PCTBACHMOR), the percentage of vacant
properties (PCTVACANT), the percentage of housing units that are
detached single family houses (PCTSINGLES), the natural logarithm of the
number of people below the poverty line (LNNBELPOV), and median
household income (MEDHHINC). The dependent variable is the natural
logarithm of median house value (LNMEDHVAL). Although OLS provides a
useful baseline for understanding these relationships, its assumption
that observations are independent can be problematic when analyzing
geographically distributed data. Neighboring census areas may share
similar socioeconomic conditions and housing-market characteristics,
resulting in spatial autocorrelation that conventional regression models
do not adequately address.

The purpose of this report is to extend the previous analysis by
applying spatial lag regression, spatial error regression, and
geographically weighted regression (GWR). These methods address spatial
relationships in different ways: spatial lag regression accounts for
dependence in the outcome across neighboring areas, spatial error
regression models spatial dependence in unobserved influences, and GWR
allows relationships between predictors and housing values to vary
geographically.

\section{Methods}\label{methods}

\subsection{Spatial Autocorrelation}\label{spatial-autocorrelation}

The 1st Law of Geography states that near things are more related than
distant things. Moran's I measures the degree of spatial autocorrelation
in our data, telling us whether data points with similar values are near
each other. The Moran's I formula produces one number that explains
whether our data is positively autocorrelated (positive values, data
with similar values are near to each other), negatively autocorrelated
(negative values, data with high values are near to values with low
values), or have no spatial autocorrelation (very close to 0 or, more
precisely -1/(n-1), the data is randomly scattered). The formula is as
follows:

\[I = \frac{n}{\sum_{i=1}^{n}\sum_{j=1}^{n} w_{ij}} \cdot \frac{\sum_{i=1}^{n}\sum_{j=1}^{n} w_{ij}(x_i - \bar{x})(x_j - \bar{x})}{\sum_{i=1}^{n}(x_i - \bar{x})^2}\]

The explanation of the terms is as follows:

\begin{itemize}
\tightlist
\item
  \(n\): the number of spatial units (for this model it is block groups)
\item
  \(x_i\): the value of the variable at location \(i\)
\item
  \(x_j\): the value of the variable at location \(j\)
\item
  \(\bar{x}\): the mean of the variable across \(n\) locations
\item
  \(w_{ij}\): the spatial weight between locations \(i\) and \(j\) taken
  from the weight matrix
\item
  \(\sum_{i}\sum_{j} w_{ij}\): the sum of all weights
\item
  \((x_i - \bar{x}) (x_j - \bar{x})\): the output of deviations from the
  mean
\item
  \(\sum_{i}(x_i - \bar{x})^2\): the sum of squared deviations from the
  mean, which shows the total variation in the variable
\end{itemize}

Throughout our report, we used a Queen contiguity weight matrix, meaning
we defined neighbors as block groups that share any point along borders.
Statisticians often use multiple spatial weight matrices in analyses to
measure the efficacy of their model. Given the various methods for
defining neighbors in spatial statistics, this gives the statistician
more tools to verify whether their model produces consistent results
with different spatial weight matrices, or whether the result is a
product of the specific spatial weight matrix they chose.

We use Moran's I as a statistical test to determine if there is spatial
autocorrelation and if so, whether the data is positively or negatively
spatially autocorrelated. Our Moran's I value was .79, indicating
positive spatial autocorrelation. Typically, negative spatial
autocorrelation ranges from -1 to near 0, and positive spatial
autocorrelation ranges from near 0 to 1. This shows that our data is
spatially clustered; in other words, data points with high values are
close to other data points with high values, and data points with low
values are near other points with low values.

We then want to know whether this pattern is statistically significant
or if it could have occurred by chance. The null hypothesis is that
there is no spatial autocorrelation, meaning there is no spatial pattern
to our data. The alternative hypothesis is that there is positive
autocorrelation. It is worth noting that the alternative hypothesis is
the presence of positive autocorrelation (not negative) because,
consistent with the First Law of Geography, we expect positive spatial
autocorrelation in geographic data and have thus decided to run a
one-sided test. To test our hypothesis, we use a permutation test, which
takes all observed values, shuffles them, randomly reassigns them to the
map, and then recalculates Moran's I. We repeat this 999 times. If fewer
than 5\% of the 999 random permuted values are as large as or larger
than the originally observed Moran's I value (when \(\alpha\) = .05), we
can reject the null hypothesis. We calculated a pseudo p-value of .001,
meaning none of the 999 permuted values were as large or larger than the
original Moran's I value. This indicates significance, and thus we can
reject the null hypothesis.

The limitation of using the global Moran's I is that it is a single
number that represents the entire data set, it does not indicate where
the spatial autocorrelation is. It tells us simply that spatial
clustering is happening. To identify more localized patterns of spatial
autocorrelation, we run Local Moran's I, also known as the Local
Indicator of Spatial Association (LISA). Instead of measuring the
overall pattern of spatial autocorrelation like we do for the global
Moran's I, the LISA goes block by block, calculating how similar that
block group is to the average of its neighbors.

Although the significance testing for the local Moran's I works
differently from the permutation test we used for the global Moran's I,
it answers the same question: is this observation unusual compared to
what a random distribution of data would produce? The following code
calculates the local Moran's I for each block group, compares each value
to its expected value under spatial randomness using a z-score, and then
derives a p-value from that z-score.

\begin{Shaded}
\begin{Highlighting}[]
\NormalTok{lisa }\OtherTok{\textless{}{-}} \FunctionTok{localmoran}\NormalTok{(philly}\SpecialCharTok{$}\NormalTok{LNMEDHVAL, queen\_weights, }\AttributeTok{zero.policy =} \ConstantTok{TRUE}\NormalTok{)}

\NormalTok{philly}\SpecialCharTok{$}\NormalTok{local\_I }\OtherTok{\textless{}{-}}\NormalTok{ lisa[, }\StringTok{"Ii"}\NormalTok{]}
\NormalTok{philly}\SpecialCharTok{$}\NormalTok{local\_p }\OtherTok{\textless{}{-}}\NormalTok{ lisa[, }\StringTok{"Pr(z != E(Ii))"}\NormalTok{]}
\end{Highlighting}
\end{Shaded}

Because the null hypothesis assumes no spatial autocorrelation, the test
calculates what Local Moran's I is expected to be if that were true. It
also calculates how much it would vary by chance alone. Then it compares
the observed Local Moran's I value to the expected value using a
z-score. This measures how many standard deviations the observed value
is from the expected value if the null hypothesis were true.

We then mapped the LISA p-values. The map shows significant p-values in
dark blue (p \textless{} .001), blue (p \textless{} .01), and light blue
(p \textless{} .05), and not significant values in pink (p ≥ .05). Pink
means that we cannot tell if that block group is part of a spatial
pattern. There is not enough evidence to say whether the block group is
similar or different from its neighbors, but the test indicates that it
could conceivably happen by chance alone.

As previously noted, neighbors in our model are defined by the Queen
matrix. Based on its relationship to the average of its neighbors, each
block group falls into one of four categories. We have found the
following schema helpful in conceptualizing these categories.

\begin{itemize}
\tightlist
\item
  \textbf{High-High:} I am a high value block group and my neighbors are
  high value block groups
\item
  \textbf{Low-Low:} I am a low value block group and my neighbors are
  low value block groups
\item
  \textbf{High-Low:} I am a high value block group and my neighbors are
  low value block groups
\item
  \textbf{Low-High:} I am a low value block group and my neighbors are
  high value block groups
\end{itemize}

We mapped these values, producing a result that appeared similar to the
LISA P-Value Map.

\subsection{OLS Regression and Assumptions
Review}\label{ols-regression-and-assumptions-review}

Ordinary Least Squares (OLS) regression is a type of linear regression
used to model the relationship between a dependent variable and one or
more independent variables. Statisticians use OLS for two primary
purposes: inference and prediction. Inferential analysis uses a sample
to draw conclusions about a population, and predictive analysis uses the
model to predict outcomes for unseen data. OLS chooses coefficients that
minimize the sum of squared residuals, defining the line or plane that
best fits the data. The OLS assumptions are as follows:

\textbf{Linearity.} The relationship between each predictor and the
dependent variable is assumed to be linear. If the true relationship
curves, a straight line will misstate the effect across parts of the
range.

\textbf{Independence of observations.} Observations, and therefore their
residuals, are assumed to be independent of one another.

\textbf{Normality of residuals.} The errors are assumed to be normally
distributed with a mean of zero. Note that this applies to the residuals
rather than to the raw variables, which need not be normal themselves.

\textbf{Homoscedasticity.} The variance of the errors is assumed to be
constant across all levels of the predicted values. If instead the
errors widen or narrow systematically, a condition called
heteroscedasticity, the standard errors and the significance tests built
on them may be unreliable.

\textbf{Multicollinearity.} Predictors can be correlated with each
other, but not so strongly (e.g., \textbar r\textbar{} \textgreater{}
0.8, or VIF above 4, and especially above 10) that the model can't
separate their effects. Severe multicollinearity makes coefficients
unstable and inflates their standard errors.

\textbf{Observations.} No fewer than 10 observations per predictor.

Our previous report goes further into depth about OLS, which you can
read here.

When the data has a spatial component, the OLS assumption that errors
are independent is often violated. We can test this by calculating the
Moran's I of the residuals, which analyzes spatial autocorrelation of
the residuals. Another way to test OLS residuals for spatial
autocorrelation is to regress them on nearby residuals. If the
independence assumption holds, block group's residuals shouldn't have
much to do with their neighbors. We visualized this by creating a
scatterplot of the OLS residuals and the weighted average of neighboring
OLS residuals. The slope of the regression line was .731, for every one
unit increase in the average residual of a block group's neighbor, that
block group's own residual is expected to increase by about .73. Simply
put, block groups whose neighbors were under or over predicted, tend to
be under or over predicted in the same direction. The first step in
calculating the slope is finding the average of the OLS residuals of
neighboring block groups. Then the OLS residuals are regressed on the
average of the neighbor's OLS residuals. The slope measures the extent
to which the two variables move together, divided by the variance of the
neighbors' average residuals. Plainly, it is the expected change of the
block group's residuals for each one-unit increase in the average
residual of its neighbor.

R, the tool we used to run the regression, can also test other
regression assumptions, such as homoscedasticity, which is also related
to the assumption of independence of errors. R has specific packages
designed to run the Breusch-Pagan, Koenker-Bassett, and White test,
which determine whether the variance of the residuals changes with the
predictors. In these tests, the null hypothesis is that the errors are
homoscedastic and the alternative hypothesis is that the errors are
heteroscedastic. A p-value \textless{} .05 means we reject the null
hypothesis, while a p-value \textgreater{} .05 means we fail to reject
the null hypothesis, indicating there is no significant evidence of
heteroscedasticity.

R can also test the assumption of normality of errors with a package
that runs the Jarque-Bera test. The null hypothesis is that the errors
are normally distributed and the alternative hypothesis is that the
errors are not normally distributed. A p-value \textless{} .05 means we
reject the null hypothesis and conclude that the residuals are not
normal, while a p-value \textgreater{} .05 means we fail to reject the
null hypothesis and conclude that the residuals are normal.

\subsection{Spatial Lag and Spatial Error
Models}\label{spatial-lag-and-spatial-error-models}

We used R to run spatial lag and spatial error regressions. Unlike OLS,
the spatial lag model assumes spatial dependence. The model accounts for
this dependence by adding an extra predictor to the OLS model: the
average dependent variable of each neighbor, or in our case, the average
house value of each block group's neighbors (W * LNMEDHVAL).

\[\begin{aligned}
\text{LNMEDHVAL} = \; & \rho W \cdot \text{LNMEDHVAL} + \beta_0 + \beta_1 \text{PCTBACHMOR} + \beta_2 \text{PCTVACANT} \\
& + \beta_3 \text{PCTSINGLES} + \beta_4 \text{LNNBELPOV} + \beta_5 \text{MEDHHINC} + \varepsilon
\end{aligned}\]

\textbf{W} is our queen weights and \textbf{LNMEDHVAL} is our dependent
variable, log median house values for all block groups.

Which means,

\(W * LNMEDHVAL\) is the spatially lagged dependent variable: the
average log median house value of each block group's neighbors.

\(\rho\) measures the strength of spatial dependence: how much a block
group's house value depends on its neighbors' house values.

\(\beta_0\) is the intercept, \(\beta_1\) through \(\beta_5\) are the
coefficients for each predictor in the error term.

We also used a spatial error model in our analysis, which accounts for
spatial dependencies in the unobserved factors that could impact our
dependent variable. In contrast to the spatial lag model that says house
value is influenced directly by surrounding house values, the spatial
error model says that house values in a similar region are subject to
the same unobserved conditions.

\[ \begin{aligned} \text{LNMEDHVAL} = \; & \beta_0 + \beta_1 \text{PCTBACHMOR} + \beta_2 \text{PCTVACANT} \\& + \beta_3 \text{PCTSINGLES} + \beta_4 \text{LNNBELPOV} + \beta_5 \text{MEDHHINC} + \varepsilon \end{aligned} \]

\[ \varepsilon = \lambda W \varepsilon + u \]

\begin{itemize}
\tightlist
\item
  \(\beta_0\) is the intercept, and \(\beta_1\) through \(\beta_5\) are
  the coefficients for each predictor.
\item
  \(\varepsilon\) is the error term from the main regression, the part
  of house values the predictors don't explain.
\item
  \(W\) is the spatially lagged error: the average error of each block
  group's neighbors, as defined by the Queen matrix.
\item
  \(\lambda\) measures how strongly a block group's error is related to
  its neighbors' errors.
\item
  \(u\) is the remaining random error, which should be independent and
  free of spatial autocorrelation.
\end{itemize}

The other OLS assumptions also pertain to Spatial Lag and Spatial Error
models, with the exception of spatial independence of observations. The
goal of these models is to account for spatial dependence, resulting in
a better fitting model with higher rates of reliability. We are able to
measure the accuracy of our model with various metrics. First, we can
look at the Moran's I of residuals. The value should be closer to 0 than
the OLS Moran's I of residuals and ideally not statistically
significant. If the spatial models work, the spatial lag term (ρ) or
spatial error term (λ) will capture the spatial patterns, instead of
leaving them in the residuals. We also compare each model by using the
Akaike Information Criterion/Schwarz Criterion, Log Likelihood, and the
Likelihood Ratio Test. These criteria are only valid when the models
being compared share the same dependent variable and are fit to the same
set of observations, which is the case for all three models in this
analysis. Each criterion also has its own characteristics and
requirements, which are discussed below.

The Akaike Information Criterion/Schwarz Criterion are measures of
goodness of fit. We use them when we want to compare models while
accounting for complexity because it penalizes each additional
parameter. When comparing AIC across models, we opt for a lower value
because it indicates the model fits well without unnecessary parameters.
AIC is only useful for comparing models. We can use the AIC to compare
nested and non-nested models.

We use Log Likelihood as a measure of fit, by scoring how well the model
matches the observed data. A higher score indicates that the observed
data is more probable under the model, meaning a better fit. We get the
score by multiplying the likelihood value of each block group. In our
data, this would require multiplying 1,720 block group values, resulting
in a number so extreme, it could easily become unmanageable. Instead,
the Log Likelihood takes the log of each value, which by the rules of
logarithms, turns the multiplication into addition. Because Log
Likelihood doesn't penalize additional parameters, it is important to
use it alongside AIC. Log Likelihood can be used when comparing nested
and non-nested models.

The Likelihood Ratio Test measures if a more complex model fits the data
significantly better than a simpler one. It uses log likelihood to
measure the fit of two models, one more complex than the other, for
example a model that accounts for spatial dependency and an OLS model.
It then compares the two log likelihood scores and runs a test to check
whether the difference is big enough to reliably not be attributable to
random chance. The Likelihood Ratio statistic follows a chi-squared
distribution where the degrees of freedom are equal to the difference in
the number of parameters between the two models. For our model, the df =
1 because each spatial model adds one parameter. The Likelihood Ratio
Test can only be used on nested models. Because of this, we cannot
compare spatial lag and spatial error models. In this case, they can
only be compared to OLS. The null hypothesis in the Likelihood Ratio
Test is always that the complex model is not a better fit than the
simpler model. Using the spatial lag and spatial error models as
examples, the null and alternative hypotheses are as follows:

For the spatial lag model:

\begin{itemize}
\tightlist
\item
  \(H_0: \rho = 0\) The spatial lag model is not a better fit than OLS.
\item
  \(H_1: \rho \neq 0\) The spatial lag model is a better fit than OLS.
\end{itemize}

For the spatial error model:

\begin{itemize}
\tightlist
\item
  \(H_0: \lambda = 0\) The spatial error model is not a better fit than
  OLS.
\item
  \(H_1: \lambda \neq 0\) The spatial error model is a better fit than
  OLS.
\end{itemize}

If the p-value is less than 0.05, we reject the null hypothesis and
conclude that the spatial model fits significantly better than OLS.

\subsection{Geographically Weighted
Regression}\label{geographically-weighted-regression}

We ran our Geographically Weighted Regression (GWR) analyses in R. GWR
runs a separate local regression for each block group, using a
proximity-based weighting scheme that gives closer neighbors more
influence on the local regression coefficient than farther neighbors.
Simpson's paradox describes the seemingly enigmatic condition in which a
global pattern in the dataset disappears or reverses when the data is
divided into subgroups. A single global regression may misrepresent
local relationships that cannot be captured with a single statistic.
This paradox is common in spatial data because relationships between the
dependent variable and predictors vary across space, a state called
spatial non-stationarity. GWR uniquely addresses this. In the case of
our model, the relationship between vacancy and house values may differ
substantially between neighborhoods in Philadelphia, but we wouldn't be
able to tell based on a single global statistic.

GWR uses what is called a bandwidth to determine which neighboring
observations are included in the local analysis. There are two types of
bandwidths: fixed and adaptive. Fixed bandwidth uses a fixed distance
from the regression point outward in all directions, this means the area
also remains constant for each local regression. In other words, a fixed
bandwidth uses the same sized net to capture its neighboring
observations and it is most effective when observations are relatively
uniformly distributed across space. Conversely, adaptive bandwidth is
best used when observations vary widely across space, when events are
concentrated in one area and sparse in others. Unlike the fixed
bandwidth's same sized net analogy, adaptive bandwidth's range is
dictated by the number of neighboring observations it captures. The
net's size may change, but the number of neighboring observations
remains the same. We use adaptive bandwidth in our model because the
observations in our data are not uniformly distributed across space,
they vary quite a bit given where in the city the block group is. Block
groups are small and dense in Center City and larger in the Northeast
and Northwest areas of Philadelphia. If we were to use fixed bandwidth
for this data, we would capture many block groups downtown and very few
at the edges of the city.

GWR runs by attributing a value to each coefficient at each location.
This means that our model ran a weighted regression for every one of the
1,720 block groups. It does this by selecting a block group, assigning
weights to nearby block groups based on the bandwidth, running a
weighted regression using those weights, saving those coefficients as
those block groups' local coefficients, and then repeating that process
for each block group. The GWR equation for our model is as follows:

\[\begin{aligned}
\text{LNMEDHVAL}_i = \; & \beta_0(u_i, v_i) + \beta_1(u_i, v_i)\,\text{PCTBACHMOR}_i + \beta_2(u_i, v_i)\,\text{PCTVACANT}_i \\
& + \beta_3(u_i, v_i)\,\text{PCTSINGLES}_i + \beta_4(u_i, v_i)\,\text{LNNBELPOV}_i + \beta_5(u_i, v_i)\,\text{MEDHHINC}_i + \varepsilon_i
\end{aligned}\]

\begin{itemize}
\tightlist
\item
  \(i\) refers to an individual block group
\item
  \((u_i,v_i)\) are the coordinates of block group \(i\)
\item
  \(\beta_0(u_i,v_i)\) is the local intercept at block group \(i\)
\item
  \(\beta_1(u_i,v_i)\) through \(\beta_5(u_i,v_i)\) are the local
  coefficients for each predictor at block group \(i\), which can vary
  across space
\item
  \(\varepsilon_i\) is the error term for block group \(i\)
\end{itemize}

The local coefficients are estimated using weighted least squares:

\[ \hat{\beta}(u_i, v_i) = (X^T W_i X)^{-1} X^T W_i y \]

The OLS assumptions are the same for GWR. One assumption even becomes
more stringent: no severe multicollinearity. This is because GWR makes
multicollinearity more likely and potentially more potent. Each local
regression relies primarily on neighboring observations, which means
limited variation and less total amount of data per regression point.
Also, because nearby block groups tend to have similar characteristics,
predictors that are only moderately correlated across the city can be
highly correlated within a local neighborhood. When predictors are
highly correlated, the regression cannot identify which variable
produces which outcome, resulting in unstable coefficients and
potentially creating spatial patterns that may look deceptively
meaningful. We can check for multicollinearity using the condition
number. The condition number measures the degree of multicollinearity in
a set of predictors. It is calculated for each local regression and if
it reads greater than 30, we generally cannot trust its results. While
stationarity isn't named as a formal assumption of OLS, it is implicit
in the mechanics of the equation. GWR also implies stationarity,
although it takes on a different meaning. OLS assumes global
stationarity; GWR assumes local stationarity.

\section{Results}\label{results}

\subsection{Spatial Autocorrelation}\label{spatial-autocorrelation-1}

\textbf{Global Moran's I}

The global Moran's I for LNMEDHVAL, calculated with the Queen contiguity
weight matrix, is 0.794. If house values were randomly spread across the
city with no spatial pattern, Moran's I would be slightly less than
zero. Our value is much higher than that, which means there is strong
positive spatial autocorrelation. In other words, block groups with high
house values tend to be next to other high-value block groups, and block
groups with low house values tend to be next to other low-value block
groups.

\begin{longtable}[]{@{}lr@{}}
\caption{Queen Contiguity Weight Matrix: Neighbor
Summary}\tabularnewline
\toprule\noalign{}
Statistic & Result \\
\midrule\noalign{}
\endfirsthead
\toprule\noalign{}
Statistic & Result \\
\midrule\noalign{}
\endhead
\bottomrule\noalign{}
\endlastfoot
Number of block groups & 1,720 \\
Total number of neighbor links & 10,526 \\
Average number of neighbors & 6.12 \\
Minimum number of neighbors & 1 \\
Maximum number of neighbors & 27 \\
Block groups with no neighbors & 0 \\
\end{longtable}

To check whether this pattern could have happened by chance, we compared
our observed Moran's I to 999 random permutations. The shuffled Moran's
I values all fall close to zero, while our observed value of 0.794 is
far to the right of all of them. None of the 999 permuted values was as
large as ours, which gives a pseudo p-value of 0.001. So we reject the
null hypothesis and conclude that LNMEDHVAL is significantly and
positively spatially autocorrelated.

\begin{longtable}[]{@{}lr@{}}
\caption{Global Moran's I: Significance Tests for Median House
Values}\tabularnewline
\toprule\noalign{}
Statistic & Result \\
\midrule\noalign{}
\endfirsthead
\toprule\noalign{}
Statistic & Result \\
\midrule\noalign{}
\endhead
\bottomrule\noalign{}
\endlastfoot
Moran's I & 0.794 \\
Expected I & -0.0006 \\
Standard deviate (z) & 56.32 \\
Analytical p-value & \textless0.001 \\
Observed rank & 1000 \\
Number of simulations & 1000 \\
Permutation pseudo p-value & 0.001 \\
Alternative hypothesis & Greater \\
\end{longtable}

\begin{figure}

{\centering \includegraphics[width=1\linewidth]{Assignment-2_files/figure-latex/fig-global-moran-hist-1} 

}

\caption{Distribution of Moran's I from 999 Permutations (observed value in red)}\label{fig:fig-global-moran-hist}
\end{figure}

\textbf{Local Moran's I}

The global Moran's I gives us one number for the whole city, but it
doesn't tell us where the clustering is happening. To find that out, we
used the Local Moran's I (LISA). LISA looks at each block group one at a
time and compares its median house value to the values of its Queen
neighbors. A block group counts as ``high'' if its value is above the
1761 block group average and ``low'' if it is below.

The significance map shows which block groups have a local Moran's I
that is statistically significant (p \textless{} 0.05) and which do not.
The significant block groups line up closely with the clusters on the
cluster map, and the block groups that are not significant fill most of
the space between them.

\begin{figure}

{\centering \includegraphics[width=1\linewidth]{Assignment-2_files/figure-latex/fig-lisa-p-1} 

}

\caption{LISA: Statistically Significant Block Groups}\label{fig:fig-lisa-p}
\end{figure}

The cluster map sorts the significant block groups into four categories.
High-High block groups have high house values and are surrounded by
neighbors with high house values. These clusters show up in three main
areas: Northwest Philadelphia (around Chestnut Hill, Mount Airy, and
Roxborough), the Far Northeast, and Center City and the neighborhoods
around it. Low-Low block groups have low house values and are surrounded
by neighbors with low house values. The largest Low-Low cluster is in
North Philadelphia, with smaller ones in West and Southwest
Philadelphia.

\begin{figure}

{\centering \includegraphics[width=1\linewidth]{Assignment-2_files/figure-latex/fig-lisa-clusters-1} 

}

\caption{LISA: Clusters of Median House Values}\label{fig:fig-lisa-clusters}
\end{figure}

High-Low and Low-High block groups are the odd ones out, where a block
group's value is very different from its neighbors. High-Low means the
block group has a high value but its neighbors are low, and Low-High
means the block group has a low value but its neighbors are high. There
are very few of these on the map. The Low-High block groups mostly sit
at the edges of the High-High clusters, and the few High-Low block
groups appear near the Low-Low clusters. This makes sense given our high
global Moran's I: in most of Philadelphia, block groups look like their
neighbors rather than different from them.

The rest of the city is not significant. This does not mean those block
groups have average house values. It means the similarity between each
block group and its neighbors is not strong enough for us to say it
isn't just chance.

\subsection{A Review of OLS Regression and
Assumptions}\label{a-review-of-ols-regression-and-assumptions}

\textbf{OLS regression results}

The OLS regression of LNMEDHVAL on PCTBACHMOR, PCTVACANT, PCTSINGLES,
LNNBELPOV, and MEDHHINC shows that four of the five predictors are
statistically significant (p \textless{} 0.001): PCTBACHMOR and MEDHHINC
are positively related to median house value, while PCTVACANT and
LNNBELPOV are negatively related. PCTSINGLES is not significant (p =
0.445), and together the predictors explain about 69\% of the variance
in LNMEDHVAL (R² = 0.6886, adjusted R² = 0.6877).

\begin{longtable}[]{@{}lrrrr@{}}
\caption{OLS Regression Results}\tabularnewline
\toprule\noalign{}
Variable & Coefficient & Std. Error & t-Statistic & P-value \\
\midrule\noalign{}
\endfirsthead
\toprule\noalign{}
Variable & Coefficient & Std. Error & t-Statistic & P-value \\
\midrule\noalign{}
\endhead
\bottomrule\noalign{}
\endlastfoot
(Intercept) & 10.6196 & 0.0607 & 174.9985 & \textless0.001 \\
PCTBACHMOR & 0.0180 & 0.0006 & 31.4127 & \textless0.001 \\
PCTVACANT & -0.0159 & 0.0010 & -16.2998 & \textless0.001 \\
PCTSINGLES & 0.0005 & 0.0007 & 0.7642 & 0.445 \\
LNNBELPOV & -0.0313 & 0.0090 & -3.4647 & \textless0.001 \\
MEDHHINC & 0.0000 & 0.0000 & 12.0395 & \textless0.001 \\
\end{longtable}

\textbf{Tests on heteroscedasticity results}

To check whether the variance of the OLS residuals is constant, we ran
three tests: the Breusch-Pagan test, the Koenker-Bassett test, and the
White test. For all three, the null hypothesis is that the residuals are
homoscedastic and the alternative is that they are heteroscedastic. The
Breusch-Pagan statistic is 196.52 (df = 5, p \textless{} 0.001), the
Koenker-Bassett statistic is 74.67 (df = 5, p \textless{} 0.001), and
the White statistic is 189.90 (df = 20, p \textless{} 0.001).

\begin{longtable}[]{@{}lrrr@{}}
\caption{Tests for Heteroscedasticity of OLS Residuals}\tabularnewline
\toprule\noalign{}
Test & Statistic & df & P-value \\
\midrule\noalign{}
\endfirsthead
\toprule\noalign{}
Test & Statistic & df & P-value \\
\midrule\noalign{}
\endhead
\bottomrule\noalign{}
\endlastfoot
Breusch-Pagan & 196.52 & 5 & \textless0.001 \\
Koenker-Bassett & 74.67 & 5 & \textless0.001 \\
White & 189.90 & 20 & \textless0.001 \\
\end{longtable}

The three tests are consistent with each other. All of them have
p-values well below 0.05, so we reject the null hypothesis in every case
and conclude that the residuals are heteroscedastic. The Koenker-Bassett
test is especially useful here because it still works when the residuals
are not normally distributed, which is the case for our model, as shown
in the next section.

\begin{figure}

{\centering \includegraphics[width=1\linewidth]{Assignment-2_files/figure-latex/fig-std-resid-fitted-1} 

}

\caption{Standardized OLS Residuals vs. Predicted Values}\label{fig:fig-std-resid-fitted}
\end{figure}

This is not fully consistent with our conclusion in Homework 1. There,
we looked at the plot of standardized residuals against predicted
values, described it as a random cloud, and concluded the residuals were
homoscedastic. Looking at the plot again for the current model, most
points do form a cloud around zero, but the spread is wider at lower
predicted values and narrower at higher ones. This pattern is easy to
miss by eye, and with 1,720 block groups the formal tests can pick up
even small differences in variance. The tests therefore give us a more
reliable answer than the plot alone.

\textbf{Tests on normality of errors (Jarque-Bera test)}

The Jarque-Bera test checks whether the OLS residuals are normally
distributed. The null hypothesis is that the residuals are normal, and
the alternative is that they are not. The Jarque-Bera statistic is
763.34 (df = 2, p \textless{} 0.001), so we reject the null hypothesis
and conclude that the residuals are not normally distributed.

\begin{longtable}[]{@{}lrrr@{}}
\caption{Jarque-Bera Test for Normality of OLS Residuals}\tabularnewline
\toprule\noalign{}
Test & Statistic & df & P-value \\
\midrule\noalign{}
\endfirsthead
\toprule\noalign{}
Test & Statistic & df & P-value \\
\midrule\noalign{}
\endhead
\bottomrule\noalign{}
\endlastfoot
Jarque-Bera & 763.34 & 2 & \textless0.001 \\
\end{longtable}

\begin{figure}

{\centering \includegraphics[width=1\linewidth]{Assignment-2_files/figure-latex/fig-std-resid-hist-1} 

}

\caption{Histogram of Standardized OLS Residuals}\label{fig:fig-std-resid-hist}
\end{figure}

This is different from our conclusion in Homework 1, where we said the
histogram of standardized residuals looked normal. At first glance the
histogram is bell-shaped and centered at zero. Looking more closely, it
is taller and narrower in the middle than a normal curve and has long
tails, with a few residuals reaching about −4 and +6. The Jarque-Bera
test looks specifically at the shape of the distribution, meaning
skewness and how heavy the tails are. With 1,720 observations, it can
detect these departures from normality even when the histogram looks
close to normal by eye.

\textbf{Scatter-plot of OLS residuals with weighted residuals}

This scatterplot shows each block group's OLS residual plotted against
the weighted residual, which is the average residual of its Queen
neighbors. If the OLS assumption of independent errors held, there would
be no pattern and the line would be flat. Instead, the points and the
fitted line slope upward. This means block groups tend to have errors in
the same direction as their neighbors. When the model overpredicts house
values in one block group, it usually overpredicts in the block groups
around it too, and the same is true for underprediction.

\begin{figure}

{\centering \includegraphics[width=1\linewidth]{Assignment-2_files/figure-latex/fig-resid-scatter-1} 

}

\caption{Scatter-plot of OLS Residuals with Weighted Residuals}\label{fig:fig-resid-scatter}
\end{figure}

\begin{longtable}[]{@{}lrrrr@{}}
\caption{Regression of OLS Residuals on Weighted
Residuals}\tabularnewline
\toprule\noalign{}
Variable & Estimate & Std. Error & t-Statistic & P-value \\
\midrule\noalign{}
\endfirsthead
\toprule\noalign{}
Variable & Estimate & Std. Error & t-Statistic & P-value \\
\midrule\noalign{}
\endhead
\bottomrule\noalign{}
\endlastfoot
(Intercept) & -0.004 & 0.007 & -0.522 & 0.601 \\
WT\_RESID & 0.731 & 0.033 & 22.000 & \textless0.001 \\
\end{longtable}

\begin{longtable}[]{@{}lr@{}}
\caption{Model Fit Statistics: OLS Residual Regression}\tabularnewline
\toprule\noalign{}
Model Statistic & Value \\
\midrule\noalign{}
\endfirsthead
\toprule\noalign{}
Model Statistic & Value \\
\midrule\noalign{}
\endhead
\bottomrule\noalign{}
\endlastfoot
Residual Standard Error & 0.3106 \\
Degrees of Freedom & 1718 \\
R-squared & 0.2198 \\
Adjusted R-squared & 0.2193 \\
F-statistic & 483.98 \\
Model p-value & \textless{} 0.001 \\
\end{longtable}

The slope b is 0.731, so for every one unit increase in the neighbors'
average residual, a block group's own residual goes up by about 0.73.
The slope is statistically significant (t = 22.00, p \textless{} 0.001),
which means there is significant spatial autocorrelation in the OLS
residuals. The neighbors' residuals alone explain about 22\% of the
variation in a block group's own residual (R² = 0.22), which is a lot
for something OLS assumes should be zero.

\textbf{Scatter-plot of Moran's I}

The Moran's I scatterplot plots each block group's OLS residual on the
x-axis against the average residual of its Queen neighbors on the
y-axis. The line through the points slopes upward, and its slope equals
the Moran's I value of 0.301. Most of the points fall in the bottom-left
and top-right parts of the plot. This means block groups where the model
overpredicts tend to be next to other block groups where it also
overpredicts, and block groups where it underpredicts tend to be next to
others where it underpredicts.

\begin{figure}

{\centering \includegraphics[width=1\linewidth]{Assignment-2_files/figure-latex/fig-moran-ols-1} 

}

\caption{Moran's I Scatterplot of OLS Residuals}\label{fig:fig-moran-ols}
\end{figure}

The Moran's I of 0.301 is far above the value we would expect if the
residuals had no spatial pattern, which is slightly less than zero. The
analytical test confirms that this is significant (z = 21.38, p
\textless{} 0.001). The OLS residuals are positively spatially
autocorrelated, even though the OLS model assumes they are independent.

\textbf{Results of 999 permutations for OLS regression residuals}

To check whether the spatial autocorrelation in the OLS residuals could
be due to chance, we compared the observed Moran's I to 999 random
permutations of the residuals. Our observed Moran's I of 0.301 ranked
1,000 out of 1,000, meaning none of the 999 permuted values was as large
as ours. This gives a pseudo p-value of 0.001, so we reject the null
hypothesis and conclude that there is significant positive spatial
autocorrelation in the OLS residuals.

\begin{longtable}[]{@{}lr@{}}
\caption{Moran's I Significance Tests: OLS Residuals}\tabularnewline
\toprule\noalign{}
Statistic & Result \\
\midrule\noalign{}
\endfirsthead
\toprule\noalign{}
Statistic & Result \\
\midrule\noalign{}
\endhead
\bottomrule\noalign{}
\endlastfoot
Moran's I & 0.301 \\
Expected I & -0.0006 \\
Standard deviate (z) & 21.38 \\
Analytical p-value & \textless0.001 \\
Observed rank & 1000 \\
Number of simulations & 1000 \\
Permutation pseudo p-value & 0.001 \\
Alternative hypothesis & Greater \\
\end{longtable}

This is a problem because OLS assumes the errors are independent of each
other. When they are not, the model is missing something about location,
and its standard errors and p-values may not be reliable. This is why we
go on to test the spatial lag and spatial error models.

Moran's I and the slope b of the weighted residuals tell the same story.
Both are positive and significant, and both show that a block group's
error is similar to its neighbors' errors. The numbers are different
(0.30 and 0.73) because they are calculated in opposite directions.
Moran's I is the slope when the neighbors' residuals are plotted against
a block group's own residual, and b is the slope when a block group's
residual is plotted against its neighbors' residuals.

\subsection{Spatial Lag and Spatial Error Regression
Results}\label{spatial-lag-and-spatial-error-regression-results}

\subsubsection{Spatial Lag Regression
Results}\label{spatial-lag-regression-results}

\textbf{Coefficients}

The spatial lag model adds one new term to the OLS model, W\_LNMEDHVAL,
which is the average log median house value of each block group's Queen
neighbors. Its coefficient, rho (ρ), is 0.624 and is highly significant
(z = 34.26, p \textless{} 0.001). This means a block group's house value
is strongly tied to the house values around it, even after accounting
for the five predictors. Because both the block group's value and its
neighbors' average are logged, we can read this roughly as: a 1\%
increase in the neighbors' average median house value is associated with
about a 0.62\% increase in the block group's own median house value,
holding the other predictors constant.

\begin{longtable}[]{@{}lrrrr@{}}
\caption{Spatial Lag Regression Results}\tabularnewline
\toprule\noalign{}
Variable & Coefficient & Std. Error & z-Statistic & P-value \\
\midrule\noalign{}
\endfirsthead
\toprule\noalign{}
Variable & Coefficient & Std. Error & z-Statistic & P-value \\
\midrule\noalign{}
\endhead
\bottomrule\noalign{}
\endlastfoot
Rho (W\_LNMEDHVAL) & 0.6242 & 0.0182 & 34.2558 & \textless0.001 \\
(Intercept) & 3.9420 & 0.1990 & 19.8075 & \textless0.001 \\
PCTBACHMOR & 0.0075 & 0.0005 & 14.4990 & \textless0.001 \\
PCTVACANT & -0.0073 & 0.0008 & -9.6436 & \textless0.001 \\
PCTSINGLES & 0.0008 & 0.0005 & 1.5377 & 0.1241 \\
LNNBELPOV & -0.0114 & 0.0068 & -1.6669 & 0.0955 \\
MEDHHINC & 0.0000 & 0.0000 & 8.0439 & \textless0.001 \\
\end{longtable}

PCTBACHMOR, PCTVACANT and MEDHHINC are still significant (p \textless{}
0.001) and have the same signs as in OLS, while PCTSINGLES is still not
significant (p = 0.124). The main change is LNNBELPOV, which was
significant in OLS (p \textless{} 0.001) but is no longer significant in
the spatial lag model (p = 0.096). All of the coefficients also get
smaller. For example, the coefficient for PCTBACHMOR drops from 0.0180
in OLS to 0.0075, and PCTVACANT drops from −0.0159 to −0.0073. This
happens because some of the effect that OLS gave to the predictors is
now explained by the neighbors' house values. In other words, OLS was
partly giving the predictors credit for the influence of location.

\textbf{Breusch-Pagan test of spatial lag residuals}

The Breusch-Pagan test shows that the spatial lag residuals are still
heteroscedastic (BP = 325.71, df = 5, p \textless{} 0.001). The
statistic is even higher than it was for OLS (196.52), so accounting for
spatial dependence did not fix the uneven size of the errors.

\begin{longtable}[]{@{}lrrr@{}}
\caption{Breusch-Pagan Test: Spatial Lag Residuals}\tabularnewline
\toprule\noalign{}
Test & Statistic & df & P-value \\
\midrule\noalign{}
\endfirsthead
\toprule\noalign{}
Test & Statistic & df & P-value \\
\midrule\noalign{}
\endhead
\bottomrule\noalign{}
\endlastfoot
Breusch-Pagan & 325.71 & 5 & \textless0.001 \\
\end{longtable}

\textbf{Comparing Spatial Lag and OLS: AIC, Schwarz Criterion, Log
Likelihood and Likelihood Ratio Test}

Compared to OLS, the spatial lag model performs better on every measure
of fit. Its AIC (462.83) and Schwarz Criterion (506.43) are much lower
than those of OLS (1297.35 and 1335.50), and lower values mean a better
fit even after penalizing the extra parameter. Its log likelihood is
also much higher (−223.42 compared to −641.68). The Likelihood Ratio
Test confirms this. The null hypothesis is that the spatial lag model
does not fit better than OLS (ρ = 0), and the LR statistic is 836.52
with 1 degree of freedom (p \textless{} 0.001). We reject the null
hypothesis and conclude that the spatial lag model fits significantly
better than OLS.

\begin{longtable}[]{@{}
  >{\raggedright\arraybackslash}p{(\linewidth - 10\tabcolsep) * \real{0.1412}}
  >{\raggedleft\arraybackslash}p{(\linewidth - 10\tabcolsep) * \real{0.1765}}
  >{\raggedleft\arraybackslash}p{(\linewidth - 10\tabcolsep) * \real{0.0941}}
  >{\raggedleft\arraybackslash}p{(\linewidth - 10\tabcolsep) * \real{0.2118}}
  >{\raggedleft\arraybackslash}p{(\linewidth - 10\tabcolsep) * \real{0.2471}}
  >{\raggedleft\arraybackslash}p{(\linewidth - 10\tabcolsep) * \real{0.1294}}@{}}
\caption{Model Comparison: OLS and Spatial Lag}\tabularnewline
\toprule\noalign{}
\begin{minipage}[b]{\linewidth}\raggedright
Model
\end{minipage} & \begin{minipage}[b]{\linewidth}\raggedleft
Log Likelihood
\end{minipage} & \begin{minipage}[b]{\linewidth}\raggedleft
AIC
\end{minipage} & \begin{minipage}[b]{\linewidth}\raggedleft
Schwarz Criterion
\end{minipage} & \begin{minipage}[b]{\linewidth}\raggedleft
LR Statistic vs.~OLS
\end{minipage} & \begin{minipage}[b]{\linewidth}\raggedleft
LR P-value
\end{minipage} \\
\midrule\noalign{}
\endfirsthead
\toprule\noalign{}
\begin{minipage}[b]{\linewidth}\raggedright
Model
\end{minipage} & \begin{minipage}[b]{\linewidth}\raggedleft
Log Likelihood
\end{minipage} & \begin{minipage}[b]{\linewidth}\raggedleft
AIC
\end{minipage} & \begin{minipage}[b]{\linewidth}\raggedleft
Schwarz Criterion
\end{minipage} & \begin{minipage}[b]{\linewidth}\raggedleft
LR Statistic vs.~OLS
\end{minipage} & \begin{minipage}[b]{\linewidth}\raggedleft
LR P-value
\end{minipage} \\
\midrule\noalign{}
\endhead
\bottomrule\noalign{}
\endlastfoot
OLS & -641.68 & 1297.35 & 1335.50 & - & - \\
Spatial Lag & -223.42 & 462.83 & 506.43 & 836.52 & \textless0.001 \\
\end{longtable}

\textbf{Moran's I scatterplot of spatial lag residuals}

The Moran's I of the spatial lag residuals is −0.072, compared to 0.301
for the OLS residuals. In the Moran's I scatterplot, the line is almost
flat and slopes slightly downward, while the OLS line clearly sloped
upward. In the 999 permutations, the observed value ranked 1 out of
1,000, giving a pseudo p-value of 0.999, so there is no positive spatial
autocorrelation left in the residuals. The value is slightly negative,
though, and the analytical test (z = −5.09) suggests this small negative
autocorrelation is not just chance. The model's Lagrange Multiplier test
for remaining residual autocorrelation is also significant (49.35, p
\textless{} 0.001). So the spatial lag model fixes most of the problem
but slightly overcorrects.

\begin{figure}

{\centering \includegraphics[width=1\linewidth]{Assignment-2_files/figure-latex/fig-moran-lag-1} 

}

\caption{Moran's I Scatterplot of Spatial Lag Residuals}\label{fig:fig-moran-lag}
\end{figure}

\begin{longtable}[]{@{}lr@{}}
\caption{Moran's I Significance Tests: Spatial Lag
Residuals}\tabularnewline
\toprule\noalign{}
Statistic & Result \\
\midrule\noalign{}
\endfirsthead
\toprule\noalign{}
Statistic & Result \\
\midrule\noalign{}
\endhead
\bottomrule\noalign{}
\endlastfoot
Moran's I & -0.072 \\
Expected I & -0.0006 \\
Standard deviate (z) & -5.09 \\
Analytical p-value & 1 \\
Observed rank & 1 \\
Number of simulations & 1000 \\
Permutation pseudo p-value & 0.999 \\
Alternative hypothesis & Greater \\
\end{longtable}

\textbf{Which model performs better}

Overall, the spatial lag model does much better than OLS. It has a lower
AIC and Schwarz Criterion, a higher log likelihood and a significant
Likelihood Ratio Test, and its residuals show far less spatial
autocorrelation. The one problem it does not solve is
heteroscedasticity.

\subsubsection{Spatial Error Regression
Results}\label{spatial-error-regression-results}

\textbf{Coefficients}

The spatial error model does not add the neighbors' house values as a
predictor. Instead, it lets the errors of neighboring block groups be
related to each other through the lambda (λ) term. Lambda is 0.803 and
is highly significant (z = 47.17, p \textless{} 0.001). This means the
part of house value that our predictors cannot explain is strongly
shared between neighbors. When the model underpredicts a block group, it
tends to underpredict its neighbors too, because they share things the
model does not measure, such as housing quality, school quality, parks
or transit access. For every one unit increase in the neighbors' average
error, a block group's own error goes up by about 0.80.

\begin{longtable}[]{@{}lrrrr@{}}
\caption{Spatial Error Regression Results}\tabularnewline
\toprule\noalign{}
Variable & Coefficient & Std. Error & z-Statistic & P-value \\
\midrule\noalign{}
\endfirsthead
\toprule\noalign{}
Variable & Coefficient & Std. Error & z-Statistic & P-value \\
\midrule\noalign{}
\endhead
\bottomrule\noalign{}
\endlastfoot
Lambda & 0.8030 & 0.0170 & 47.1723 & \textless{} 0.001 \\
(Intercept) & 10.6822 & 0.0590 & 180.9761 & \textless{} 0.001 \\
PCTBACHMOR & 0.0087 & 0.0007 & 11.8359 & \textless{} 0.001 \\
PCTVACANT & -0.0056 & 0.0009 & -6.3389 & \textless{} 0.001 \\
PCTSINGLES & 0.0018 & 0.0006 & 2.8346 & 0.00459 \\
LNNBELPOV & -0.0154 & 0.0074 & -2.0728 & 0.03820 \\
MEDHHINC & 0.0000 & 0.0000 & 7.7644 & \textless{} 0.001 \\
\end{longtable}

All five predictors are significant in the spatial error model.
PCTBACHMOR, PCTVACANT and MEDHHINC are significant (p \textless{} 0.001)
with the same signs as in OLS. LNNBELPOV is also still significant,
though less strongly than in OLS (p = 0.038 compared to p \textless{}
0.001). The biggest change is PCTSINGLES, which was not significant in
OLS (p = 0.445) but becomes significant in the spatial error model (p =
0.005), and its coefficient triples from 0.0005 to 0.0018. Most of the
other coefficients get smaller than in OLS, though not by as much as in
the spatial lag model. This suggests that once shared unmeasured
neighborhood factors are accounted for, the share of single-family homes
does matter for house values.

\textbf{Breusch-Pagan test of spatial error residuals}

The Breusch-Pagan test shows that the spatial error residuals are still
heteroscedastic (BP = 57.87, df = 5, p \textless{} 0.001). The statistic
is much lower than for OLS (196.52) and the spatial lag model (325.71),
though. So the spatial error model reduces the problem of uneven errors
but does not get rid of it.

\begin{longtable}[]{@{}lrrr@{}}
\caption{Breusch-Pagan Test: Spatial Error Residuals}\tabularnewline
\toprule\noalign{}
Test & Statistic & df & P-value \\
\midrule\noalign{}
\endfirsthead
\toprule\noalign{}
Test & Statistic & df & P-value \\
\midrule\noalign{}
\endhead
\bottomrule\noalign{}
\endlastfoot
Breusch-Pagan & 57.87 & 5 & \textless0.001 \\
\end{longtable}

\textbf{Comparing Spatial Error and OLS: AIC, Schwarz Criterion, Log
Likelihood and Likelihood Ratio Test}

Compared to OLS, the spatial error model also performs better on every
measure of fit. Its AIC (702.65) and Schwarz Criterion (746.26) are much
lower than those of OLS (1297.35 and 1335.50), and its log likelihood is
much higher (−343.33 compared to −641.68). The Likelihood Ratio Test
confirms this. The null hypothesis is that the spatial error model does
not fit better than OLS (λ = 0), and the LR statistic is 596.70 with 1
degree of freedom (p \textless{} 0.001). We reject the null hypothesis
and conclude that the spatial error model fits significantly better than
OLS.

\begin{longtable}[]{@{}
  >{\raggedright\arraybackslash}p{(\linewidth - 10\tabcolsep) * \real{0.1609}}
  >{\raggedleft\arraybackslash}p{(\linewidth - 10\tabcolsep) * \real{0.1724}}
  >{\raggedleft\arraybackslash}p{(\linewidth - 10\tabcolsep) * \real{0.0920}}
  >{\raggedleft\arraybackslash}p{(\linewidth - 10\tabcolsep) * \real{0.2069}}
  >{\raggedleft\arraybackslash}p{(\linewidth - 10\tabcolsep) * \real{0.2414}}
  >{\raggedleft\arraybackslash}p{(\linewidth - 10\tabcolsep) * \real{0.1264}}@{}}
\caption{Model Comparison: OLS and Spatial Error}\tabularnewline
\toprule\noalign{}
\begin{minipage}[b]{\linewidth}\raggedright
Model
\end{minipage} & \begin{minipage}[b]{\linewidth}\raggedleft
Log Likelihood
\end{minipage} & \begin{minipage}[b]{\linewidth}\raggedleft
AIC
\end{minipage} & \begin{minipage}[b]{\linewidth}\raggedleft
Schwarz Criterion
\end{minipage} & \begin{minipage}[b]{\linewidth}\raggedleft
LR Statistic vs.~OLS
\end{minipage} & \begin{minipage}[b]{\linewidth}\raggedleft
LR P-value
\end{minipage} \\
\midrule\noalign{}
\endfirsthead
\toprule\noalign{}
\begin{minipage}[b]{\linewidth}\raggedright
Model
\end{minipage} & \begin{minipage}[b]{\linewidth}\raggedleft
Log Likelihood
\end{minipage} & \begin{minipage}[b]{\linewidth}\raggedleft
AIC
\end{minipage} & \begin{minipage}[b]{\linewidth}\raggedleft
Schwarz Criterion
\end{minipage} & \begin{minipage}[b]{\linewidth}\raggedleft
LR Statistic vs.~OLS
\end{minipage} & \begin{minipage}[b]{\linewidth}\raggedleft
LR P-value
\end{minipage} \\
\midrule\noalign{}
\endhead
\bottomrule\noalign{}
\endlastfoot
OLS & -641.68 & 1297.35 & 1335.50 & - & - \\
Spatial Error & -343.33 & 702.65 & 746.26 & 596.7 & \textless0.001 \\
\end{longtable}

\textbf{Moran's I scatterplot of spatial error residuals}

The Moran's I of the spatial error residuals is −0.094, compared to
0.301 for the OLS residuals. In the Moran's I scatterplot, the line is
almost flat and slopes slightly downward, while the OLS line clearly
sloped upward. In the 999 permutations, the observed value ranked 1 out
of 1,000, giving a pseudo p-value of 0.999, so there is no positive
spatial autocorrelation left in the residuals. Like the spatial lag
model, though, the value is slightly negative, and the analytical test
(z = −6.63) suggests this small negative autocorrelation is not just
chance. So the spatial error model also fixes most of the problem but
slightly overcorrects.

\begin{figure}

{\centering \includegraphics[width=1\linewidth]{Assignment-2_files/figure-latex/fig-moran-error-1} 

}

\caption{Moran's I Scatterplot of Spatial Error Residuals}\label{fig:fig-moran-error}
\end{figure}

\begin{longtable}[]{@{}lr@{}}
\caption{Moran's I Significance Tests: Spatial Error
Residuals}\tabularnewline
\toprule\noalign{}
Statistic & Result \\
\midrule\noalign{}
\endfirsthead
\toprule\noalign{}
Statistic & Result \\
\midrule\noalign{}
\endhead
\bottomrule\noalign{}
\endlastfoot
Moran's I & -0.094 \\
Expected I & -0.0006 \\
Standard deviate (z) & -6.63 \\
Analytical p-value & 1 \\
Observed rank & 1 \\
Number of simulations & 1000 \\
Permutation pseudo p-value & 0.999 \\
Alternative hypothesis & Greater \\
\end{longtable}

\textbf{Which model performs better}

Overall, the spatial error model does much better than OLS. It has a
lower AIC and Schwarz Criterion, a higher log likelihood and a
significant Likelihood Ratio Test, and its residuals show far less
spatial autocorrelation. It also reduces heteroscedasticity, although it
does not remove it.

\subsubsection{Comparing Spatial Lag and Spatial
Error}\label{comparing-spatial-lag-and-spatial-error}

\textbf{AIC and Schwarz Criterion}

We cannot use the Likelihood Ratio Test to compare the spatial lag and
spatial error models, because they are not nested. Neither model is a
simpler version of the other. We can still compare them using the AIC
and Schwarz Criterion, which work for non-nested models. The spatial lag
model has the lower AIC (462.83 compared to 702.65) and the lower
Schwarz Criterion (506.43 compared to 746.26), so it is the
better-fitting model.

\begin{longtable}[]{@{}
  >{\raggedright\arraybackslash}p{(\linewidth - 10\tabcolsep) * \real{0.1609}}
  >{\raggedleft\arraybackslash}p{(\linewidth - 10\tabcolsep) * \real{0.1724}}
  >{\raggedleft\arraybackslash}p{(\linewidth - 10\tabcolsep) * \real{0.0920}}
  >{\raggedleft\arraybackslash}p{(\linewidth - 10\tabcolsep) * \real{0.2069}}
  >{\raggedleft\arraybackslash}p{(\linewidth - 10\tabcolsep) * \real{0.2414}}
  >{\raggedleft\arraybackslash}p{(\linewidth - 10\tabcolsep) * \real{0.1264}}@{}}
\caption{Model Comparison: OLS, Spatial Lag and Spatial
Error}\tabularnewline
\toprule\noalign{}
\begin{minipage}[b]{\linewidth}\raggedright
Model
\end{minipage} & \begin{minipage}[b]{\linewidth}\raggedleft
Log Likelihood
\end{minipage} & \begin{minipage}[b]{\linewidth}\raggedleft
AIC
\end{minipage} & \begin{minipage}[b]{\linewidth}\raggedleft
Schwarz Criterion
\end{minipage} & \begin{minipage}[b]{\linewidth}\raggedleft
LR Statistic vs.~OLS
\end{minipage} & \begin{minipage}[b]{\linewidth}\raggedleft
LR P-value
\end{minipage} \\
\midrule\noalign{}
\endfirsthead
\toprule\noalign{}
\begin{minipage}[b]{\linewidth}\raggedright
Model
\end{minipage} & \begin{minipage}[b]{\linewidth}\raggedleft
Log Likelihood
\end{minipage} & \begin{minipage}[b]{\linewidth}\raggedleft
AIC
\end{minipage} & \begin{minipage}[b]{\linewidth}\raggedleft
Schwarz Criterion
\end{minipage} & \begin{minipage}[b]{\linewidth}\raggedleft
LR Statistic vs.~OLS
\end{minipage} & \begin{minipage}[b]{\linewidth}\raggedleft
LR P-value
\end{minipage} \\
\midrule\noalign{}
\endhead
\bottomrule\noalign{}
\endlastfoot
OLS & -641.68 & 1297.35 & 1335.50 & - & - \\
Spatial Lag & -223.42 & 462.83 & 506.43 & 836.52 & \textless0.001 \\
Spatial Error & -343.33 & 702.65 & 746.26 & 596.70 & \textless0.001 \\
\end{longtable}

\textbf{Which model performs better}

The spatial lag model also has a higher log likelihood (−223.42 compared
to −343.33), and its residual Moran's I (−0.072) is closer to zero than
that of the spatial error model (−0.094). The one area where the spatial
error model does better is heteroscedasticity. Its Breusch-Pagan
statistic (57.87) is much lower than that of the spatial lag model
(325.71), although the residuals of both models are still
heteroscedastic.

Both models are a large improvement over OLS, but the spatial lag model
performs best overall. This suggests that house values in Philadelphia
are directly shaped by the values of neighboring block groups, and not
only by unmeasured factors that neighbors happen to share.

\subsubsection{Geographically Weighted Regression
Results}\label{geographically-weighted-regression-results}

\textbf{Global GWR results}

The GWR model uses an adaptive bisquare kernel with a bandwidth of 107
nearest neighbors, which was chosen by minimizing the AICc. This means
each local regression is based on the 107 block groups closest to it,
with closer block groups given more weight.

The GWR model explains much more of the variance in LNMEDHVAL than OLS.
Its overall R-squared is 0.851 (adjusted R-squared = 0.819), compared to
0.689 for OLS (adjusted R-squared = 0.688). By letting the relationships
between the predictors and house values change across the city, GWR
explains about 16 percentage points more of the variation in house
values than a single citywide regression.

\begin{longtable}[]{@{}lr@{}}
\caption{Global GWR Diagnostics}\tabularnewline
\toprule\noalign{}
Statistic & Value \\
\midrule\noalign{}
\endfirsthead
\toprule\noalign{}
Statistic & Value \\
\midrule\noalign{}
\endhead
\bottomrule\noalign{}
\endlastfoot
RSS.gw & 101.364 \\
AIC & 245.944 \\
AICc & 557.927 \\
enp & 309.599 \\
edf & 1410.401 \\
gw.R2 & 0.851 \\
gwR2.adj & 0.819 \\
BIC & 40.001 \\
\end{longtable}

GWR also has the lowest AIC of all four models. Its AIC is 245.94,
compared to 462.83 for the spatial lag model, 702.65 for the spatial
error model and 1297.35 for OLS. Based on the AIC, GWR fits the data
best. This improvement comes at a cost in complexity, though. Because
GWR runs a separate regression for every block group, it uses about 310
effective parameters, compared to 7 or 8 in the other models. Its AICc,
which penalizes extra parameters more heavily, is 557.93, which is
higher than the AIC of the spatial lag model.

\textbf{Moran's I scatterplot of GWR residuals}

The Moran's I of the GWR residuals is 0.030, compared to 0.301 for OLS,
−0.072 for the spatial lag model and −0.094 for the spatial error model.
In the Moran's I scatterplot, the line is almost flat, while the OLS
line clearly sloped upward. The GWR residuals have the least spatial
autocorrelation of all four models, since 0.030 is the closest value to
zero. The small amount that remains is still positive and statistically
significant, however. In the 999 permutations, the observed value ranked
987 out of 1,000, giving a pseudo p-value of 0.013. So GWR removes most
of the spatial autocorrelation in the OLS residuals, but neighboring
block groups still have slightly similar errors.

\begin{figure}

{\centering \includegraphics[width=1\linewidth]{Assignment-2_files/figure-latex/fig-moran-gwr-1} 

}

\caption{Moran's I Scatterplot of GWR Residuals}\label{fig:fig-moran-gwr}
\end{figure}

\begin{longtable}[]{@{}lr@{}}
\caption{Moran's I Significance Tests: GWR Residuals}\tabularnewline
\toprule\noalign{}
Statistic & Result \\
\midrule\noalign{}
\endfirsthead
\toprule\noalign{}
Statistic & Result \\
\midrule\noalign{}
\endhead
\bottomrule\noalign{}
\endlastfoot
Moran's I & 0.030 \\
Expected I & -0.0006 \\
Standard deviate (z) & 2.20 \\
Analytical p-value & 0.0139 \\
Observed rank & 987 \\
Number of simulations & 1000 \\
Permutation pseudo p-value & 0.013 \\
Alternative hypothesis & Greater \\
\end{longtable}

\textbf{Local regression results: coefficient-to-standard-error maps}

To see where each predictor has a possibly significant relationship with
house values, we divided each local coefficient by its standard error. A
ratio above 2 (dark blue) suggests a significant positive relationship,
a ratio below −2 (dark red) suggests a significant negative
relationship, and ratios between −2 and 2 (pink and light blue) are not
significant.

\textbf{Bachelor's degree or higher (PCTBACHMOR).} The percentage of
residents with a bachelor's degree or higher has a significant positive
relationship with house values across most of the city, including
Northwest, West, Southwest and South Philadelphia, Center City and parts
of the Northeast. The relationship is positive but not significant in a
band running through North Philadelphia and the lower and Far Northeast.
This is the most consistent predictor in the model, which matches its
strong significance in the global models.

\begin{figure}

{\centering \includegraphics[width=1\linewidth]{Assignment-2_files/figure-latex/fig-gwr-pctbachmor-1} 

}

\caption{Local Coefficient-to-SE Ratio: Bachelor's Degree}\label{fig:fig-gwr-pctbachmor}
\end{figure}

\textbf{Vacant housing units (PCTVACANT).} The percentage of vacant
housing units is negative almost everywhere, but it is significant only
in some parts of the city. The strongest negative relationships are in
West and Southwest Philadelphia, parts of South Philadelphia and the
Germantown area in the Northwest. In much of the Northeast and the outer
Northwest, the relationship is negative but not significant. This
suggests vacancy drags down house values most in neighborhoods where it
is already concentrated.

\begin{figure}

{\centering \includegraphics[width=1\linewidth]{Assignment-2_files/figure-latex/fig-gwr-pctvacant-1} 

}

\caption{Local Coefficient-to-SE Ratio: Percentage of Vacant Households}\label{fig:fig-gwr-pctvacant}
\end{figure}

\textbf{Single-family homes (PCTSINGLES).} The percentage of
single-family homes shows the clearest change from place to place. It
has a significant positive relationship with house values in the Far
Northeast and the outer Northwest, where detached single-family homes
are common. In parts of North, West and South Philadelphia and along the
Delaware River, the relationship is negative and in some pockets
significantly negative. This helps explain why PCTSINGLES was not
significant in OLS. Its positive and negative effects in different parts
of the city cancel each other out when only one citywide coefficient is
estimated.

\begin{figure}

{\centering \includegraphics[width=1\linewidth]{Assignment-2_files/figure-latex/fig-gwr-pctsingles-1} 

}

\caption{Local Coefficient-to-SE Ratio: Percentage of Single-Family Homes}\label{fig:fig-gwr-pctsingles}
\end{figure}

\textbf{Households in poverty (LNNBELPOV).} The log of the number of
households in poverty is negative or close to zero in most of the city,
but it is significant in only a few pockets, mainly in North
Philadelphia, the Germantown area, near the Delaware River and in South
Philadelphia. This matches the global results, where LNNBELPOV became
weaker or lost significance once spatial dependence was accounted for.

\begin{figure}

{\centering \includegraphics[width=1\linewidth]{Assignment-2_files/figure-latex/fig-gwr-lnnbelpov-1} 

}

\caption{Local Coefficient-to-SE Ratio: Log of Households in Poverty}\label{fig:fig-gwr-lnnbelpov}
\end{figure}

\textbf{Median household income (MEDHHINC).} Median household income has
a positive relationship with house values in most of the city, but it is
significant only in some areas. The strongest positive relationships are
in the outer Northwest, a large area of the lower Northeast, and
scattered parts of West, Center City and South Philadelphia. In much of
the Far Northeast and North Philadelphia, the relationship is positive
but not significant, and in some areas it is slightly negative but not
significant. One area in the Northwest, around Germantown, shows a
significant negative relationship, meaning higher incomes are linked to
lower house values there. This may reflect local conditions the model
does not capture, such as older or larger housing that sells for less,
or neighborhoods going through change. Overall, income matters for house
values in many parts of Philadelphia, but not equally everywhere.

\begin{figure}

{\centering \includegraphics[width=1\linewidth]{Assignment-2_files/figure-latex/fig-gwr-medhhinc-1} 

}

\caption{Local Coefficient-to-SE Ratio: Median Household Income}\label{fig:fig-gwr-medhhinc}
\end{figure}

\textbf{Local R-squared}

The local R-squared map shows how well the GWR model explains house
values in each part of the city. Values range from about 0.4 to above
0.8. The model performs best (local R-squared above 0.8) in the
Northwest, the Far Northeast and the area around West Philadelphia and
Center City. It performs worst (local R-squared around 0.4 to 0.6) in
North Philadelphia, along the Delaware River and at the southern tip of
South Philadelphia. In these areas, house values are likely shaped by
factors the model does not include, such as housing condition,
redevelopment pressure or industrial land uses.

\begin{figure}

{\centering \includegraphics[width=1\linewidth]{Assignment-2_files/figure-latex/fig-gwr-local-r2-1} 

}

\caption{Local R-squared from Geographically Weighted Regression}\label{fig:fig-gwr-local-r2}
\end{figure}

Overall, the map shows that a single citywide model does not fit every
part of Philadelphia equally well, which is exactly what GWR is designed
to show.

\section{Discussion}\label{discussion}

This report examined spatial variation in median housing values across
Philadelphia block groups using three regression methods: Spatial Lag,
Spatial Error, and Geographically Weighted Regression (GWR). The results
of these regression models were compared with each other and with the
Ordinary Least Squares (OLS) model from the previous report. The
dependent variable was the log-transformed median housing value, while
the independent variables were percentage of residents with a bachelor's
degree or higher, percentage of vacant households, percentage of
single-family homes, and percentage of the population under the poverty
line.

The first step in the analysis was creating a Queen weight matrix to
determine Moran's I. The Moran's I (assumed global) tells us how
spatially autocorrelated the entire model is, 1 being positively
spatially autocorrelated, -1 being negatively spatially autocorrelated,
and 0 being not spatially autocorrelated. We then used a Monte Carlo
test to check if the Moran's I values are accurately representative of
the model, and plotted a histogram of all 999 permutations. Following
the global Moran's I calculations, we then calculated the local Moran's
I, also known as LISA. Using LISA, we were able to calculate a Moran's I
for each individual block group, thus allowing us to know how spatially
autocorrelated each block group is. We then created a map of these
results. Since we decided to complete our analysis in R rather than
GeoDa, our OLS regression from the previous report is already calculated
and saved into our data frame.

We then calculated the weights (i.e.~spatially lagged) residuals from
the OLS regression. An important distinction is the difference between
spatially lagged residuals and residuals from a Spatial Lag model. A
spatially lagged residual is a weighted average of neighboring
observations' residuals, calculated by applying a spatial weights matrix
to the residual vector. It describes how the errors estimated for nearby
areas relate to one another and is often used to investigate residual
spatial autocorrelation. By contrast, a Spatial Lag model includes a
spatially lagged dependent variable as an explanatory component,
representing the relationship between an area's outcome and the outcomes
of neighboring areas. Its residuals are the differences between observed
and fitted values after accounting for that spatial-lag term. Thus, the
spatial lag in the model refers to neighboring outcomes, whereas
spatially lagged residuals refer to neighboring model errors. The two
concepts are related to spatial dependence but are not interchangeable.
Distinguishing them is essential when interpreting model specifications
and evaluating whether spatial structure remains unexplained after
estimation.

We then calculated the Moran's I for the OLS residuals and mapped the
results. We repeated these steps for Spatial Lag, Spatial Error, and
Geographically Weighted Regression models. Comparing these methods
provided a way to evaluate both the relationships between socioeconomic
characteristics and housing values and the extent to which spatial
dependence influenced model performance.

It is important to note that each model operates within specific
limitations and under specific assumptions. OLS provides a baseline
estimate of the relationships between the independent variables and
housing values, but it assumes that observations are independent and
that the model errors are not spatially autocorrelated. These
assumptions may be problematic in housing markets, where neighboring
areas often share economic conditions, amenities, development patterns,
and historical influences. The Spatial Lag model accounts for the
influence of neighboring observations through a spatially lagged
dependent variable, whereas the Spatial Error model accounts for spatial
dependence in the unobserved factors represented by the error term. GWR
extends the analysis by allowing relationships between predictors and
housing values to vary geographically, revealing the possibility that
socioeconomic characteristics have different associations with housing
values across Philadelphia. Additionally, Queen-contiguity weights
define neighbors according to shared boundaries, which may not fully
represent housing-market relationships across physical or socioeconomic
boundaries. Finally, the analysis identifies associations rather than
causal relationships: omitted characteristics, such as housing quality,
accessibility, environmental conditions, and historical patterns of
investment, may also influence housing values. The best-performing
method should be identified by comparing model fit, information
criteria, and residual spatial autocorrelation, rather than assuming
that a more complex model is necessarily superior. In particular, lower
AIC or AICc values and weaker residual spatial autocorrelation can
support a model's selection, although these measures must be interpreted
alongside the models' assumptions and comparability. All of this
considered, we decided that the best regression model for this scenario
is Spatial Lag.

Ultimately, understanding the spatial distribution of housing values
requires attention not only to socioeconomic predictors but also to the
geographic relationships between neighborhoods. By incorporating spatial
dependence and local variation, this analysis provides a more
geographically nuanced framework for investigating the factors
associated with residential property values in Philadelphia.

\end{document}
